\section{Arithmetic input and exact convergence domains}
\label{sec:convergence}

\subsection{The common-base identities}

If \(a,b\) are multiplicatively dependent, take \(a^u=b^v\) with
\(\gcd(u,v)=1\). For every prime \(p\),
\(u\,v_p(a)=v\,v_p(b)\), so \(v\mid v_p(a)\) and \(u\mid v_p(b)\).
Thus \(a=c^v,b=c^u\) for the integer
\[
c=\prod_p p^{v_p(a)/v}.
\]
The coprime exponent pair and then \(c\) are unique by the same
valuation comparison. The integer \(c\) need not be power-free.
Set \(r=v,s=u\).

For positive \(u\geq v\), subtraction gives
\[
(c^u-1)-c^{u-v}(c^v-1)=c^{u-v}-1.
\]
The Euclidean algorithm on the exponents therefore proves
\begin{equation}\label{eq:same-base}
\gcd(c^u-1,c^v-1)=c^{\gcd(u,v)}-1.
\end{equation}
In particular, the dependent coefficient is
\(C_{a,b}(n,m)=c^{\gcd(rn,sm)}-1\).

\subsection{Multiplicatively independent bases}

We use the following consequence of the general \(S\)-unit theorem
of Corvaja--Zannier \citep[Corollary 1 and (1.3), public version 2]{cz2005}.
For each finite prime set \(S\) and each \(\epsilon>0\), positive
integer \(S\)-units \(u,v\) that are multiplicatively independent
satisfy
\[
\gcd(u-1,v-1)<\max(u,v)^\epsilon
\]
apart from finitely many exceptional pairs.
This is a statement about two independent exponents, not merely a
diagonal estimate with a shared exponent.

Take \(S\) to contain the prime divisors of \(ab\). If \(a,b\) are
multiplicatively independent, so are \(a^n,b^m\) for every \(n,m\geq1\).
Absorbing the finite exceptional set into a constant \(K_\epsilon\)
gives the uniform estimate
\begin{equation}\label{eq:cz-bound}
C_{a,b}(n,m)
\leq K_\epsilon
\exp\bigl(\epsilon\max(n\log a,m\log b)\bigr).
\end{equation}
Let \(0<u,v<1\) and \(L=\max(\log a,\log b)\).
Choose \(\epsilon>0\) such that
\(u e^{\epsilon L}<1\) and \(v e^{\epsilon L}<1\).
Since the maximum in \eqref{eq:cz-bound} is at most \(L(n+m)\),
the double series is bounded on \(|x|\leq u,|y|\leq v\) by a
constant times
\[
\sum_{n,m\geq1}(u e^{\epsilon L})^n(v e^{\epsilon L})^m<\infty.
\]
This proves locally normal absolute convergence on \(\D^2\).

For \(y\ne0\) and \(|x|\geq1\), a fixed row \(m\geq1\) has divergent
absolute sum, since \(C_{a,b}(n,m)\geq1\). The symmetric assertion
holds for \(x\ne0\), \(|y|\geq1\).
Although the series vanishes pointwise on either coordinate axis,
an axis point outside \(\D^2\) has no neighborhood of absolute
convergence. This proves the first part of
Theorem~\ref{thm:classification}.

\subsection{Multiplicatively dependent bases}

Put \(u=|x|,v=|y|\), initially with \(0<u,v<1\).
The inequality in \eqref{eq:omega} is equivalent to
\(c u^{1/r}v^{1/s}<1\). Choose \(0<\theta<1\) sufficiently close to
one that \(c u^{\theta/r}v^{\theta/s}<1\).
For \(g=\gcd(rn,sm)\), the bounds \(n\geq g/r\) and \(m\geq g/s\)
imply
\begin{align*}
C_{a,b}(n,m)u^nv^m
&\leq c^g u^nv^m\\
&\leq \bigl(c u^{\theta/r}v^{\theta/s}\bigr)^g
u^{(1-\theta)n}v^{(1-\theta)m}\\
&\leq u^{(1-\theta)n}v^{(1-\theta)m}.
\end{align*}
The last expression is summable. The strict inequalities allow the
same estimate on a slightly larger closed polydisc around each
nonaxis point of \(\Omega_{r,s,c}\).
At an axis point in this domain, choose a small positive radius for
the zero coordinate and a radius less than one for the other so that
the same strict inequality holds. Thus the series is locally normally
convergent throughout \(\Omega_{r,s,c}\).

Conversely, if \(c^{rs}u^sv^r\geq1\), the terms on the distinguished
primitive ray \((n,m)=(sk,rk)\) are
\[
(c^{rsk}-1)u^{sk}v^{rk},
\]
and they do not tend to zero. If instead \(u\geq1\) or \(v\geq1\)
with both coordinates nonzero, the fixed-row argument applies.
Together with the axis convention, these facts give exactly
\eqref{eq:omega} as the open absolute-convergence domain.
They establish convergence geometry; they do not preclude
meromorphic passage through parts of its boundary.
